\documentclass[12pt]{article}
\usepackage[T1]{fontenc}
\usepackage[utf8]{inputenc}
\usepackage{lmodern}
\usepackage{textcomp}
\usepackage{enumitem}
\usepackage{geometry}
\geometry{margin=1in}
\begin{document}
\begin{center}
Answer Key - Sociology - Undergraduate Level Assessment 20 \
\small (Answers are ordered exactly as in the question paper.)
\end{center}

\section*{Multiple Choice}
\begin{enumerate}[label=\arabic*.]
\item \textbf{Answer:} A
\\ \textit{Options:} A) we live in a world of superficial, fragmented images; B) no theory is better than any other: 'anything goes'; C) society has changed and we need new kinds of theory; D) all of the above
\\ \textit{Note:} Post-modernist writers contend that contemporary society is characterized by a proliferation of superficial and fragmented images, reflecting the disintegration of coherent narratives and the rise of simulacra in culture and media.
\item \textbf{Answer:} D
\\ \textit{Options:} A) the ruling class, or bourgeoisie, who exploit the proletariat; B) a capitalist class dependent on property ownership and advantaged life chances; C) an alignment of classes with shared interests but no state power; D) a state elite whose members are drawn overwhelmingly from a power bloc
\\ \textit{Note:} The term 'power elite' as introduced by Scott (1991) refers to a distinct group within the state that holds significant power and influence, composed predominantly of individuals from a dominant social bloc, thereby highlighting the concentration of power among a select few in society.
\item \textbf{Answer:} D
\\ \textit{Options:} A) a complex network of interaction at a micro-level; B) a source of conflict, inequality, and alienation; C) an unstable structure of social relations; D) a normative framework of roles and institutions
\\ \textit{Note:} Structural-Functionalists view society as a complex system of interrelated roles and institutions that work together to maintain social order and stability, emphasizing the importance of norms in guiding behavior and the functioning of the social structure.
\item \textbf{Answer:} D
\\ \textit{Options:} A) Prime Minister, Archbishop of York, Viscounts of England; B) Marquesses of England, Earls of Great Britain, King's Brothers; C) Esquires, Serjeants of Law, Dukes' Eldest Sons; D) King's Grandsons, Lord High Treasurer, Companions of the Bath
\\ \textit{Note:} The correct answer reflects the traditional hierarchy of British nobility, where the king's grandsons hold the highest status as royal family members, followed by the Lord High Treasurer, a significant political figure, and then the Companions of the Bath, who are typically lesser-ranking honors within the British honors system.
\item \textbf{Answer:} B
\\ \textit{Options:} A) the proportion of people living alone has fallen to 29\%; B) many people are cohabiting in long term relationships; C) the upward curve of remarriages compensates for the drop in first marriages; D) all of the above
\\ \textit{Note:} The correct answer highlights that the rise in cohabitation reflects a shift in societal attitudes towards commitment and partnership, leading many to choose long-term living arrangements without formal marriage, thereby contributing to the perceived decline of traditional marriage.
\end{enumerate}

\section*{True / False}
\begin{enumerate}[label=\arabic*.]
\setcounter{enumi}{5}
\item \textbf{Answer:} False
\\ \textit{Options:} A) True; B) False
\\ \textit{Note:} Thatcher's government implemented policies that prioritized reducing state welfare spending and promoting self-reliance, which resulted in decreased financial support for single parents, students, and the unemployed.
\item \textbf{Answer:} True
\\ \textit{Options:} A) True; B) False
\\ \textit{Note:} Sociology employs diverse research methods and approaches across different contexts, allowing for a richer and more nuanced understanding of complex social phenomena.
\item \textbf{Answer:} True
\\ \textit{Options:} A) True; B) False
\\ \textit{Note:} Comte's term 'positivism' indeed refers to an approach in sociology that emphasizes the importance of empirical evidence and scientific methods in understanding social phenomena, asserting that knowledge should be derived from observable facts rather than speculation or metaphysics.
\item \textbf{Answer:} True
\\ \textit{Options:} A) True; B) False
\\ \textit{Note:} Pilcher's identification of soap operas as a 'feminine genre' reflects the dual portrayal of women in these narratives, showcasing their roles as both caretakers within domestic spaces and as individuals asserting their independence, thus highlighting the complexities of female identity in contemporary society.
\item \textbf{Answer:} False
\\ \textit{Options:} A) True; B) False
\\ \textit{Note:} The market model of state welfare typically emphasizes targeted assistance and benefits based on individual income and needs rather than providing universal benefits to all individuals.
\end{enumerate}

\section*{Long Form Response}
\begin{enumerate}[label=\arabic*.]
\setcounter{enumi}{10}
\item \textbf{Answer:} The New Labour government, which came to power in 1997, implemented a series of educational reforms aimed at improving standards across the UK. A significant aspect of their approach involved targeted support for failing Local Education Authorities (LEAs). Rather than imposing punitive measures, the government emphasized collaboration and assistance to help these LEAs enhance their performance. This supportive strategy included increased funding, the introduction of performance indicators, and the establishment of partnerships with successful schools to share best practices. By focusing on uplifting struggling LEAs, New Labour sought to create a more equitable educational landscape, ultimately aiming to raise overall educational attainment and reduce disparities among different regions.
\\ \textit{Note:} correct because it accurately outlines New Labour's educational reform policies focused on improving failing LEAs through supportive measures rather than punitive actions, highlighting the use of increased funding and performance indicators as key strategies for enhancing educational standards.
\item \textbf{Answer:} The Conservative government of 1979 implemented several strategies aimed at diminishing the influence of the labour movement, with one of the most significant being the introduction of legislation that made all strike action illegal. This move effectively undermined the power of trade unions, which had historically been a strong advocate for workers' rights and collective bargaining. By criminalizing strikes, the government sought to deter workers from organizing and protesting for better conditions, thereby shifting the balance of power away from unions and towards employers. The implications of this approach were profound, as it not only restricted workers' rights to collective action but also led to a culture of fear and compliance within the workforce, diminishing the overall effectiveness of labor movements in advocating for social justice and fair labor practices.
\\ \textit{Note:} correct because it identifies the introduction of legislation criminalizing strike action as a key strategy by the Conservative government of 1979 to weaken trade unions and diminish their influence on workers' rights, thereby illustrating the direct impact of government policy on labor organization and collective bargaining.
\end{enumerate}

\vfill
\noindent\textit{End of Answer Key}
\end{document}